% =====================================================================
%  Seance 2 -- B. Preparer les donnees (frames 10--18)
% =====================================================================
\section{Préparer les données : doublons, manquants, aberrants}

\begin{frame}[fragile]{Étape 0 : trier et dédoublonner}
\begin{codebox}[Le rituel des 3 premières lignes de tout pipeline]
\begin{lstlisting}
b = pd.read_csv("data/nosys_metriques_brutes.csv", parse_dates=["horodatage"])

# 1. trier : la collecte arrive dans le desordre
b = b.sort_values(["serveur", "horodatage"]).reset_index(drop=True)

# 2. dedoublonner : 200 lignes exactement identiques (re-tentative reseau)
print(b.duplicated().sum())        # -> 200
b = b.drop_duplicates()

# 3. verifier la completude : nb de points attendus par serveur
b.groupby("serveur").size()
\end{lstlisting}
\end{codebox}
\begin{columns}[t]
\begin{column}{0.48\textwidth}
\begin{ebox}[Ce que révèle l'étape 3 sur le fichier brut]
7 serveurs à $\approx$ 8 600 points, \textbf{srv-07 à 2 880} : la
granularité 15 min se voit immédiatement dans le compte.
\end{ebox}
\end{column}
\begin{column}{0.48\textwidth}
\begin{wbox}[Pourquoi trier AVANT tout le reste]
Diff, fenêtres glissantes et interpolation supposent l'ordre
temporel. Sur un fichier « brûlé » (ordre mélangé), une moyenne
mobile calculée avant le tri est \textbf{mathématiquement fausse}
--- sans qu'aucune erreur ne soit levée.
\end{wbox}
\end{column}
\end{columns}
\end{frame}

\begin{frame}{Valeurs manquantes : trois mécanismes, trois réponses}
\begin{center}
\footnotesize
\begin{tabular}{@{}p{2.4cm}p{4.8cm}p{5.0cm}@{}}
\toprule
\textbf{Mécanisme} & \textbf{Ce que c'est} & \textbf{Réponse adaptée}\\
\midrule
\textbf{MCAR} & manquant \emph{au hasard} (capteur flottant) &
imputation simple acceptable\\
\addlinespace[1pt]
\textbf{MAR} & manquant \emph{lié à une autre variable observée} &
modèle d'imputation (KNN, régression)\\
\addlinespace[1pt]
\textbf{MNAR} & manquant \emph{pour une raison liée à sa valeur}
(le CPU à 100 \% fait planter l'agent !) &
ne jamais imputer aveuglément : documenter, traiter le trou comme signal\\
\bottomrule
\end{tabular}
\end{center}
\vspace{1mm}
\begin{columns}[t]
\begin{column}{0.52\textwidth}
\begin{ebox}[Dans le fichier brut NOSYS]
\textbf{5 508} cellules manquantes (1,73 \% des cellules numériques),
réparties uniformément sur les 5 colonnes (1 072 à 1 123 chacune) :
profil \textbf{MCAR} --- capteurs flottants. Imputation justifiée,
\textbf{après} vérification du profil.
\end{ebox}
\end{column}
\begin{column}{0.44\textwidth}
\begin{kbox}[Le test à faire systématiquement]
Comparer le taux de manquants \textbf{pendant les incidents} et
\textbf{hors incidents}. Si les manquants se concentrent pendant les
incidents $\Rightarrow$ MNAR : l'absence de donnée \emph{est} une information.
\end{kbox}
\end{column}
\end{columns}
\end{frame}

\begin{frame}{Imputer : forward-fill contre interpolation}
\begin{center}
\includegraphics[width=0.94\textwidth]{f13_imputation.png}
\end{center}
\begin{columns}[t]
\begin{column}{0.52\textwidth}
\begin{ebox}[Erreur mesurée sur CPU (600 points masqués, srv-01)]
\begin{itemize}
\item \textbf{Forward-fill} (recopier la dernière valeur) :
erreur absolue moyenne \textbf{1,018} points de CPU.
\item \textbf{Interpolation linéaire} :
\textbf{0,743} --- 27 \% mieux, car le signal évolue en tendance.
\end{itemize}
\end{ebox}
\end{column}
\begin{column}{0.44\textwidth}
\begin{wbox}[La règle du pas]
Jamais d'imputation au-delà d'1--2 pas. Au-delà (trou de 40 min),
interpoler \textbf{invente} des données : mieux vaut marquer le trou
et l'exclure des fenêtres d'entraînement.
\end{wbox}
\end{column}
\end{columns}
\end{frame}

\begin{frame}{Valeurs aberrantes : z-score et IQR en pratique}
\begin{center}
\includegraphics[width=0.92\textwidth]{f12_outliers.png}
\end{center}
\begin{columns}[t]
\begin{column}{0.5\textwidth}
\begin{ebox}[Sur la latence du fichier brut]
Après exclusion des 12 pics capteur (latence $>$ 400 ms) :
\begin{itemize}
\item bornes IQR : \py{[Q1 - 1.5*IQR ; Q3 + 1.5*IQR]} ;
\item bornes z-score : moyenne $\pm$ 3 écart-types.
\item Les deux méthodes se recouvrent sur une loi proche de la normale ---
l'IQR est \textbf{robuste} même si la distribution est asymétrique.
\end{itemize}
\end{ebox}
\end{column}
\begin{column}{0.46\textwidth}
\begin{wbox}[Le défaut à ne JAMAIS faire]
Supprimer tout ce qui dépasse ! Sur nos données, \textbf{les incidents
sont des valeurs extrêmes légitimes} (CPU 96 \% pendant l'incident).
\begin{itemize}
\item Distinguer \textbf{erreur de mesure} (température 1 575 °C,
CPU 350 \% : physiquement impossible) et \textbf{événement réel} (CPU 96 \%).
\item Règle métier d'abord : bornes physiques du système, puis statistiques.
\end{itemize}
\end{wbox}
\end{column}
\end{columns}
\end{frame}

\begin{frame}[fragile]{Rééchantillonner : du pas irrégulier à la grille commune}
\begin{codebox}[Normaliser la granularité (srv-07 en 15 min $\to$ 5 min)]
\begin{lstlisting}
def normaliser(g):                      # g = un serveur
    g = g.set_index("horodatage").drop(columns=["serveur"])
    # grille reguliere de 5 min sur la plage du serveur
    grille = pd.date_range(g.index.min(), g.index.max(), freq="5min")
    return g.reindex(grille)            # les trous deviennent des NaN

b = b.groupby("serveur", group_keys=False).apply(normaliser)
\end{lstlisting}
\end{codebox}
\begin{columns}[t]
\begin{column}{0.5\textwidth}
\begin{center}
\includegraphics[width=\textwidth]{f09_reechantillonnage.png}
\end{center}
\end{column}
\begin{column}{0.46\textwidth}
\begin{kbox}[Ce que fait \py{reindex}]
Il crée la grille 5 min \textbf{explicite} : les trous (agent redémarré)
apparaissent comme \py{NaN} au lieu de disparaître silencieusement ---
on peut alors les compter, les tracer, les imputer ou les exclure
en connaissance de cause.
\end{kbox}
\end{column}
\end{columns}
\end{frame}

\begin{frame}{Quiz 2 --- préparation (5 min)}
\begin{qbox}[3 questions]
\begin{enumerate}
\item 5 508 manquants uniformément répartis sur les colonnes : quel mécanisme ?\\
(a) MNAR \quad (b) MAR \quad (c) MCAR
\item Un CPU affiché à 350 \%. Que fait-on ?\\
(a) on l'impute \quad (b) c'est une erreur de mesure (physiquement impossible)
$\Rightarrow$ on la corrige en NaN puis on impute \quad (c) on le garde, c'est
un incident
\item Pourquoi l'IQR est-il souvent préféré au z-score sur des métriques d'infra ?\\
(a) il est plus rapide \quad (b) il est robuste aux distributions
asymétriques (latence, queues longues) \quad (c) il détecte plus d'outliers
\end{enumerate}
\end{qbox}
\vspace{1mm}
\begin{solbox}[Réponses]
\textbf{1 : c} \quad \textbf{2 : b} --- CPU $>$ 100 \% n'existe pas ;
CPU 96 \% pendant un incident est, lui, légitime. \quad \textbf{3 : b}
\end{solbox}
\end{frame}

\begin{frame}{À retenir de la préparation}
\begin{kbox}[Les 5 règles du pipeline de préparation]
\begin{enumerate}
\item \textbf{Trier} par (serveur, temps) avant toute fenêtre ou interpolation.
\item \textbf{Dédoublonner} puis \textbf{compter} : la granularité se voit
dans les comptes (srv-07 : 2 880 points).
\item \textbf{Réindexer} sur une grille explicite : les trous deviennent
des NaN \emph{visibles}, jamais des silences.
\item Classer les manquants (\textbf{MCAR/MAR/MNAR}) avant de choisir
l'imputation --- et ne jamais imputer au-delà d'1--2 pas.
\item Distinguer \textbf{erreur de mesure} (impossible physiquement)
et \textbf{événement réel} : les incidents SONT des valeurs extrêmes légitimes.
\end{enumerate}
\end{kbox}
\end{frame}
